在应用中, 宜用余弦级数还是正弦级数展开, 通常取决于边界条件; 见 Chapter~\ref{chap:introduction} 中的例子, 以及以下关于 Fourier 级数逐点收敛的说明. Corollary~\ref{cor:fourier-complex-hilbert-basis} 提供的正弦与余弦展开通常用于周期函数.

知道 Fourier 级数在 $L^2$ 中收敛固然很好, 但仍留下若干问题. 例如:\begin{enumerate}
\item 在什么条件下 Fourier 级数按逐点意义表示一个函数?
\item Fourier 级数能否逐项求导? 当然可以在分布意义下这样做, 但也希望知道求导后的级数是否在 $L^2$ 中收敛.
\end{enumerate}

由测度论知道, 在 $L^2$ 中收敛的序列具有几乎处处收敛的子列. 对 Fourier 级数, 实际上不必取子列; 这是直到 1966 年才被证明的一个困难定理. 一个初等得多的观察是, 若 $\sum_{n=0}^\infty|a_n|$ 收敛, 则 $\sum_{n=0}^\infty a_n\cos(n\pi x)$ 在 $[0,1]$ 上一致收敛; 利用 $\ell^2$ 中的 Cauchy-Schwarz 不等式可知, 若对任意 $\alpha>1$ 有 $\sum_{n=1}^\infty|a_n|^2n^\alpha$ 收敛, 就属于这一情形 (令 $|a_n|=(|a_n|n^{\alpha/2})(n^{-\alpha/2})$; 若 $\alpha>1$, 则序列 $n^{-\alpha/2}$ 属于 $\ell^2$).

现在设 $f\in L^2(0,1)$, 且 $f$ 的导数 (在分布意义下) 也属于 $L^2(0,1)$ (稍后将在 Sobolev 空间部分 (Chapter~\ref{chap:sobolev-spaces}) 详细研究这样的函数). 那么 $f'$ 可以展开成正弦级数或余弦级数:
\begin{equation}\label{eq:fourier-derivative-sine-cosine-expansions}
\begin{aligned}
f'(x)&=\sum_{n=1}^\infty a_n\sin(n\pi x),\\
f'(x)&=\sum_{n=0}^\infty b_n\cos(n\pi x).
\end{aligned}
\end{equation}
积分可得
\begin{equation}\label{eq:fourier-antiderivative-cosine-sine-expansions}
\begin{aligned}
f(x)&=\sum_{n=1}^\infty-\dfrac{a_n}{n\pi}\cos(n\pi x)+\alpha,\\
f(x)&=b_0x+\sum_{n=1}^\infty\dfrac{b_n}{n\pi}\sin(n\pi x)+\beta.
\end{aligned}
\end{equation}


这些表达式中的第一个表示 $f$ 的余弦级数, 且由于 \[\sum_{n=1}^\infty|a_n^2|<\infty\] 有 \[\sum_{n=1}^\infty|a_n|/n\pi<\infty,\] 即 \eqref{eq:fourier-antiderivative-cosine-sine-expansions} 中的第一个级数一致收敛. 因而, 任何具有 $L^2(0,1)$ 中导数的 $f\in L^2(0,1)$ 都具有一致收敛的余弦级数. (特别地, 这意味着任何这样的 $f$ 都连续. 这是 Sobolev 嵌入定理的一个特例.) 此外, 在 $L^2$ 收敛意义下, 该级数可以逐项求导.

另一方面, \eqref{eq:fourier-antiderivative-cosine-sine-expansions} 中的第二个级数只有在 $\beta=b_0=0$ 时才是正弦级数. 容易看出 $\beta=f(0)$ 且 $b_0=\int_0^1 f'(x)\,\mathrm{d}x=f(1)$. 因而, 任何满足 $f'\in L^2(0,1)$ 且另有 $f(0)=f(1)=0$ 的函数 $f\in L^2(0,1)$ 都具有绝对收敛的正弦级数. 这表明 Fourier 级数的收敛行为不仅受函数光滑性的影响, 还受其边界行为的影响.
